% !TeX program = pdflatex
\documentclass[
amberg,
toc=fancy,
footlinebox,
nometaot
]{OTHAWbeamer}

\usepackage[T1]{fontenc}
\usepackage[utf8]{inputenc}
\usepackage[ngerman]{babel}
\usepackage{csquotes}
\usepackage{booktabs}
\usepackage{hyperref}
\usepackage{tabularx}
\usepackage[backend=biber,style=authoryear]{biblatex}
\addbibresource{quellen.bib}

\title[PMA]{PMA}
\subtitle{Social-Matching-MVP als Hochschulprojekt}
\author{Arturo Mendez \and Team PMA}
\institute{OTH Amberg-Weiden}
\date{\today}

\begin{document}

\maketitle

\begin{frame}{Inhaltsverzeichnis}
  \tableofcontents
\end{frame}

\section{Projektkontext}

\begin{frame}{Projektziel}
  \begin{block}{Zielbild}
    PMA soll als lauffaehiger MVP zeigen, dass Social Discovery nicht nur ueber reines Personen-Matching gedacht werden kann, sondern ueber einen zusammenhaengenden Produktfluss aus Profil, Discovery, Match, Chat, Events und Kartenbezug.
  \end{block}

  \vspace{0.3cm}

  \begin{itemize}
    \item Fokus auf ein technisch vorfuehrbares End-to-End-System
    \item kein Anspruch auf vollstaendigen Produktionsumfang
    \item iterative Umsetzung ueber vier Sprints
  \end{itemize}
\end{frame}

\begin{frame}{Warum ist PMA ein Projekt?}
  \begin{columns}[T]
    \column{0.48\textwidth}
    \textbf{Projektcharakter}
    \begin{itemize}
      \item neues Produktziel
      \item begrenzter Zeitrahmen
      \item eigene Rollen und Zuständigkeiten
      \item sichtbare Risiken bei Scope und Integration
    \end{itemize}

    \column{0.48\textwidth}
    \textbf{Nicht Routine}
    \begin{itemize}
      \item kein standardisierter Ablauf
      \item Produktbild musste erst geschärft werden
      \item Anforderungen veraenderten sich durch Feedback
      \item Frontend, Backend und Datenmodell mussten gemeinsam wachsen
    \end{itemize}
  \end{columns}
\end{frame}

\section{Produkt}

\begin{frame}{Produktidee}
  \begin{itemize}
    \item PMA ist keine reine Dating-App, sondern ein Social-Discovery-MVP.
    \item Der Live-Flow fokussiert aktuell Personen-Discovery.
    \item Zusaetzliche soziale Kontexte werden ueber Interessen, Events und Kartenbezug sichtbar angebunden.
  \end{itemize}

  \vspace{0.3cm}

  \begin{block}{Kernfrage}
    Wie laesst sich Kennenlernen digital nicht nur ueber Profile, sondern auch ueber gemeinsame Interessen, spontane Treffen und Orte denken?
  \end{block}
\end{frame}

\begin{frame}{Umgesetzter MVP-Umfang}
  \begin{columns}[T]
    \column{0.48\textwidth}
    \textbf{Nutzerfluss}
    \begin{itemize}
      \item Registrierung und Login
      \item Profilpflege mit Galerie
      \item Personen-Discovery mit Swipe und Match
      \item Chat zwischen Matches
    \end{itemize}

    \column{0.48\textwidth}
    \textbf{Sozialer Kontext}
    \begin{itemize}
      \item Events erstellen, filtern und bearbeiten
      \item Kartenansicht mit Event-Markern
      \item Nutzerzonen in pseudoanonymisierter Form
      \item Orts- und Stadtvorschlaege ueber Nominatim/OSM \cite{nominatim}
    \end{itemize}
  \end{columns}
\end{frame}

\begin{frame}{Technischer Aufbau}
  \begin{tabularx}{\textwidth}{@{}>{\bfseries}l X@{}}
    Frontend & React, React Router, Tailwind, React Leaflet \\
    Backend & Node.js, Express, MongoDB, Mongoose \\
    Auth & tokenbasierte Session-Authentifizierung \\
    Tests & \texttt{node --test} mit \texttt{mongodb-memory-server} \\
    Demo & Seed-Script fuer reproduzierbare Datenlage \\
  \end{tabularx}

  \vspace{0.4cm}

  \begin{block}{Architekturidee}
    Statt isolierter Screens wurde ab Sprint 2 bewusst auf vertikale Inkremente gesetzt: UI, API und Datenmodell pro Kernfunktion gemeinsam.
  \end{block}
\end{frame}

\section{Vorgehen}

\begin{frame}{Vorgehensmodell}
  \begin{itemize}
    \item Organisation als Scrum-basiertes, iterativ-inkrementelles Softwareprojekt \cite{scrumguide}
    \item zweiwoechige Sprints mit Planning, Daily, Review und Retrospektive
    \item Product Backlog als priorisierte Arbeitsgrundlage
    \item Reviews dienten nicht nur zur Demo, sondern zur aktiven Scope-Korrektur
  \end{itemize}
\end{frame}

\begin{frame}{Rollen und Stakeholder}
  \begin{columns}[T]
    \column{0.48\textwidth}
    \textbf{Teamrollen}
    \begin{itemize}
      \item Product Owner
      \item Scrum Master
      \item Developers
    \end{itemize}

    \column{0.48\textwidth}
    \textbf{Stakeholder}
    \begin{itemize}
      \item betreuender Dozent im Modulkontext
      \item Nutzergruppe: Studierende und junge Erwachsene
      \item Feedback floss ueber Sprint Reviews in die Priorisierung ein
    \end{itemize}
  \end{columns}
\end{frame}

\begin{frame}{Meilensteine}
  \begin{tabularx}{\textwidth}{@{}l X@{}}
    M1 & Produktbild und klickbarer Frontend-Rahmen sichtbar \\
    M2 & Auth, Profil und Personen-Discovery als erster vertikaler Kernflow \\
    M3 & Matchliste, Chat sowie Event-/Map-Bereiche als erlebbare Produktteile \\
    M4 & Events und Karte an echte Backend-Daten angebunden, Demo und Doku abgabefaehig \\
  \end{tabularx}
\end{frame}

\section{Sprintverlauf}

\begin{frame}{Sprintverlauf in Kurzform}
  \begin{enumerate}
    \item \textbf{Sprint 1:} klickbares Produktbild, aber noch stark mockup-lastig
    \item \textbf{Sprint 2:} echter Kernflow mit Auth, Profil, Discovery, Swipe und Match
    \item \textbf{Sprint 3:} Match-Nutzen durch Chat, dazu erste Event- und Kartenbereiche
    \item \textbf{Sprint 4:} Backend-Anbindung fuer Events und Karte, Demo-Seed, Doku und Stabilisierung
  \end{enumerate}
\end{frame}

\begin{frame}{Wichtigstes Review-Learning}
  \begin{block}{Rueckmeldung nach Sprint 1}
    Ein reiner Frontend-Prototyp reicht nicht aus. Das Produkt musste frueh als echter Ende-zu-Ende-Flow erfahrbar werden.
  \end{block}

  \begin{itemize}
    \item Wechsel von Mockup-Denken zu vertikaler Umsetzung
    \item Personen-Discovery blieb Kern, wurde aber in ein breiteres Social-Discovery-Bild eingebettet
    \item Events und Locations wurden nicht als Deko, sondern als Produktkontext mitgedacht
  \end{itemize}
\end{frame}

\section{Einordnung}

\begin{frame}{Bewusste MVP-Grenzen}
  \begin{itemize}
    \item Bilder werden aktuell clientseitig als komprimierte Data-URLs gespeichert
    \item kein Echtzeit-Chat per WebSocket
    \item keine eigenen Swipe-Decks fuer Events oder Orte
    \item Kartenzonen sind fuer die Demo bewusst pseudoanonymisiert
  \end{itemize}
\end{frame}

\begin{frame}{Fazit}
  \begin{itemize}
    \item PMA zeigt einen technisch durchgaengigen Social-Matching-MVP.
    \item Der groesste methodische Fortschritt war der Wechsel von klickbaren Screens zu vertikalen Inkrementen.
    \item Scrum war fuer das Projekt vor allem als Struktur fuer Priorisierung, Review-Feedback und Steuerung nuetzlich.
  \end{itemize}

  \vspace{0.3cm}

  \begin{block}{Kurz gesagt}
    Aus einer fruehen UI-Idee wurde ein vorfuehrbares Produkt mit nachvollziehbarem Sprintverlauf und klaren MVP-Grenzen.
  \end{block}
\end{frame}

\appendix
\nocite{scrumguide,agilemanifesto,nominatim}

\begin{frame}[allowframebreaks]{Quellen}
  \printbibliography[heading=none]
\end{frame}

\end{document}
